\def\stat{kalenov}





\def\tit{НЕКОТОРЫЕ ОЦЕНКИ ПОИСКОВЫХ ВОЗМОЖНОСТЕЙ  


ИНТЕРНЕТ-КАТАЛОГОВ РОССИЙСКИХ БИБЛИОТЕК}





\def\titkol{Некоторые оценки поисковых возможностей  


интернет-каталогов российских библиотек}








\def\aut{Н.\,Е.~Калёнов$^1$}





\def\autkol{Н.\,Е.~Калёнов}





\titel{\tit}{\aut}{\autkol}{\titkol}





\index{Калёнов Н.\,Е.}


\index{Kalenov N.\,E.}





%{\renewcommand{\thefootnote}{\fnsymbol{footnote}}


%\footnotetext[1]


%{Работа выполнена при поддержке РФФИ (проект 15-07-02053).}}





\renewcommand{\thefootnote}{\arabic{footnote}}


\footnotetext[1]{Библиотека по естественным наукам Российской академии наук, 


\mbox{nek@benran.ru}}








 \label{st\stat}





                  \thispagestyle{headings}


                  


  \vspace*{-12pt}





  


  


   


  


  \Abst{Рассматриваются вопросы организации поиска информации в~биб\-лио\-теч\-ных 


каталогах как в~<<историческом>> аспекте (применительно к~традиционным карточным 


и~печатным каталогам), так и~в~современном, обусловленном развитием 


информационных технологий, Интернета и~универсальных поисковых машин. Проведен 


анализ и~выявлены недостатки существующих ин\-тер\-нет-ка\-та\-ло\-гов ряда российских 


научных библиотек. Рассмотрен ряд примеров некорректной обработки запросов 


каталогами ведущих биб\-лио\-тек России. Сделан вывод о~необходимости нового подхода 


к~проектированию ин\-тер\-нет-ка\-та\-ло\-гов научных биб\-лио\-тек, которые должны выдавать 


пользователю всю ту и~только ту информацию, которая его интересует. При создании 


программных оболочек ин\-тер\-нет-ка\-та\-ло\-гов необходимо, с~одной стороны, отойти 


от традиционных библиотечных представлений о~каталогах, существенно 


ограничивающих их поисковые воз\-мож\-ности, с~другой стороны~--- от логики работы 


универсальных поисковых машин Интернета, выдающих пользователю тысячи процентов 


информационного шума.}


   


  \KW{электронные каталоги; Интернет; качество информационного поиска; 


библиографические описания; сетевые технологии; поисковые запросы}





\DOI{10.14357\!/\!08696527180319} 








\vskip 10pt plus 9pt minus 6pt








      \thispagestyle{headings}


      


      \vspace*{-10pt} 


  


  Одним из основных атрибутов любой биб\-лио\-те\-ки являются каталоги, 


отражающие ее фонды. До наступления <<эры компьютеризации>> 


в~биб\-лио\-те\-ках формировались в~основном карточные, реже печатные 


каталоги. Преимущества карточных каталогов для посетителей библиотек 


очевидны~--- карточные каталоги легко актуализировать, до\-бав\-ляя и~изымая 


карточки по мере поступления или исключения из фондов конкретных 


изданий, поэтому (при чет\-ко действующем биб\-лио\-теч\-ном персонале) 


карточные каталоги в~любой момент времени отражают состояние фондов 


биб\-лио\-те\-ки. Печатные каталоги не обеспечивали оперативность отражения 


фондов, однако они могли рассылаться по почте заинтересованным 


пользователям, которые, не приходя в~биб\-лио\-те\-ку, могли получать 


информацию о~ее фондах. В~этой связи печатные каталоги обычно отражали 


либо часть фондов, относящуюся к~определенному событию (например, 


содержали перечень изданий, пред\-став\-лен\-ных на тематической или 


персональной вы\-став\-ке), либо содержали перечень изданий, поступивших за 


определенный промежуток времени. 


  


  Традиционные карточные каталоги подразделяются на два вида~--- 


алфавитные и~сис\-те\-ма\-ти\-че\-ские (существуют так\-же предметные каталоги, но 


принцип их построения не отличается от систематических). Пер\-вые служат 


для поиска конкретного издания, про которое известны фамилия автора или 


название; вторые позволяют найти издания по определенной тематике. 


Информация об издании выводится на карточку в~виде стандартизированного 


библиографического описания~[1]. Карточки в~алфавитных каталогах 


упорядочиваются по заголовкам библиографических описаний. В~заголовки 


<<основных>> описаний выводятся фамилии авторов либо названия изданий 


(для сборников, не имеющих индивидуальных авторов). Наряду с~основными 


карточками для обеспечения более полного поиска в~каталогах на одно 


издание могут формироваться добавочные карточки, в~заголовки которых 


выводятся фамилии дополнительных авторов, редакторов, со\-ста\-ви\-те\-лей 


и~т.\,п. Кроме карточек с~библиографическими описаниями в~алфавитных 


каталогах могут присутствовать также <<ссылочные>> и~<<справочные>> 


карточки. Первые используются, в~част\-ности, при наличии у~авторов 


псевдонимов и~дают ссылки на используемые в~алфавитном ряду данного 


каталога варианты имени автора (например, <<Клеменс Сэмуэл см.\ Марк 


Твен>>); вторые раскрывают аббревиатуры названий изданий или 


организаций (например, <<МТТ~--- Механика твердого тела>> или 


<<ВИНИТИ~--- Всероссийский институт научной и~технической 


информации>>). Карточки для сис\-те\-ма\-ти\-че\-ских каталогов име\-ют своими 


заголовками индексы той или иной сис\-те\-мы классификации (УДК, ББК 


и~пр.), принятой в~данной биб\-лио\-те\-ке; в~заголовки карточек предметных 


каталогов выносятся наименования предметных руб\-рик.


  


  Развитие вычислительной техники и~сетевых технологий обусловило 


возможности формирования в~библиотеках электронных каталогов. Сегодня 


трудно найти научную биб\-лио\-те\-ку, которая бы не имела своего электронного 


каталога, пред\-став\-лен\-но\-го в~Интернете. Однако в~подавляющем большинстве 


случаев эти каталоги (особенно на этапах их <<становления>>) строились 


исходя из принципов, заложенных в~традиционных карточных каталогах, 


и~были ориентированы на поиск изданий по ограниченному чис\-лу элементов 


библиографических описаний, соответственно, не использовались в~полной 


мере воз\-мож\-ности автоматизированных ин\-фор\-ма\-ци\-он\-но-по\-ис\-ко\-вых сис\-тем. 


В~качестве примера подобного подхода мож\-но привести электронный каталог ГПНТБ России 


<<WEB ИРБИС>>. Пользователю, обратившемуся к~каталогу ({\sf 


http://www.gpntb.ru/}, переход по ссылке <<Электронный каталог>>) 


предлагается ввес\-ти в~строку интересующий его термин, вы\-брав 


предварительно <<об\-ласть поиска>> (варианты: <<клю\-че\-вые слова>>, 


<<Автор>>, <<Заглавие>>, <<Год издания>>) (рис.~1).


  





\begin{figure} %fig1


\vspace*{1pt}


 \begin{center}


 \mbox{%


 \epsfxsize=128mm 


 \epsfbox{kal-1.eps}


 }


\end{center}


\vspace*{-9pt}


\Caption{Поисковая страница электронного каталога ГПНТБ России}


\end{figure}





\begin{figure} %fig2


\vspace*{1pt}


 \begin{center}


 \mbox{%


 \epsfxsize=128mm 


 \epsfbox{kal-2.eps}


 }


\end{center}


\vspace*{-9pt}


\Caption{Поисковый запрос в~электронном каталоге ГПНТБ России}


\end{figure}


  


  Предположим, пользователя интересуют издания, связанные с~протоколом 


z39.50.


  


  Если выбрать область поиска <<Ключевые слова>>, сис\-те\-ма сообщает:


  


  <<По Вашему запросу: 


$\langle$.$\rangle$\textbf{K\,=\;\,Z39}\$$\langle$.$\rangle$ ничего не найдено, 


уточните запрос>>.


  


  Если искать в~области поиска <<Заглавие>>, система выдает две записи, 


заглавие которых начинается с~z39.50 (рис.~2). 


  


  По запросу <<Жижимов>> в~об\-ласти поиска <<Автор>> сис\-те\-ма 


выдает~17~документов, первый же из которых~--- <<Жижимов,~О.\,Л. 


Построение распределенных информационных сис\-тем на основе протокола 


Z39.50 [Текст]: автореф. дис.\ \ldots\ д-ра техн. наук: 05.25.05~/ 


О.\,Л.~Жижимов.~--- Новосибирск, 2004.~--- 31~с.: ил.~--- Библиогр.:  


с.~25--31 (37~назв.)>>~--- не выдается при поиске ни по ключевым словам, 


ни по заглавию.


  


  Проявлением <<библиотечного подхода>> к~организации электронного 


каталога ГПНТБ является также текст <<доп.\ точки до\-сту\-па>> в~ниж\-ней 


час\-ти записи, относящейся к~изданиям, и~вывод значения, по которому 


возможен поиск этой записи (см.\ обе записи на рис.~2).


  


  К сожалению, WEB-ИРБИС является одной из самых распространенных 


в~стране программных оболочек, поддерживающих  


ин\-тер\-нет-ка\-та\-ло\-ги биб\-лио\-тек. Автор неоднократно указывал 


в~профессиональной биб\-лио\-теч\-ной печати на недостатки и~ошибки этих 


каталогов~[2--4], однако многие академические биб\-лио\-те\-ки продолжают их 


поддерживать. В~результате пользователи, обращаясь к~таким крупнейшим 


академическим биб\-лио\-те\-кам, как ГПНТБ СО РАН, ЦНБ УрО РАН, БАН, не 


всегда могут получить полную и~достоверную информацию о~наличии в~их 


фондах тех или иных изданий. Например, на запрос <<Экологические 


угрозы>> в~об\-ласти поиска <<Заглавие>> каталоги всех этих биб\-лио\-тек 


выдают сообщение:


  


  <<По Вашему запросу: ``($\langle$.$\rangle$\textbf{T\,=\;\,Экологические 


угрозы}\$$\langle$.$\rangle$)'' (на естественном языке: 


\textbf{ЗАГЛАВИЕ:} ``Экологические угрозы'') ничего не найдено. 


Убедитесь в~корректности запроса>>.


  


  В то же время во всех каталогах имеется книга, в~заглавие которой входит 


фрагмент <<Экологические угрозы>>: <<Тетельмин, Владимир 


Владимирович. Сланцевые углеводороды. Технологии добычи. 


Экологические угрозы: учеб. пособие для вузов~/ В.\,В.~Тетельмин, 


В.\,А.~Язев, А.\,А.~Соловьянов.~--- Долгопрудный : Интеллект, 2014.~--- 


295~с.>>.


  


  Справедливости ради надо отметить, что при вводе запроса 


<<Экологические угрозы>> в~область поиска <<ключевые слова>> эта книга 


выдается каталогами ГПНТБ СО РАН и~ЦНБ УрО РАН, но ее нет 


среди~14~книг, выдаваемых на этот запрос каталогом БАН. 


  


  Случаи, когда электронный каталог сообщает об отсутствии той или иной 


книги в~библиотеке, несмотря на то что она есть, недопустимы. При работе 


с~каталогом в~<<традиционной>> ситуации, когда читатель приходит 


в~библиотеку, такие случаи почти исключены, поскольку в~поиске нужной 


книги читателю помогают находящиеся рядом библиотечные специалисты. 


С~ин\-тер\-нет-ка\-та\-ло\-гом пользователь остается один на один, и~задача 


разработчиков~--- предоставить ему сервис, который в~максимальной 


степени должен исключить возможность выдачи неверных сведений, 


поэтому к~разработке ин\-тер\-нет-ка\-та\-ло\-гов необходимо подходить 


с~учетом этих требований. 


  


  В последнее время многие библиотеки стали впадать в~другую  


крайность~---вмес\-то расширения возможностей традиционного поиска 


пошли по пути копирования пользовательского интерфейса, принятого 


в~поисковых машинах общего назначения (Yandex, Google, Rambler и~т.\,п.), 


предоставляя для ввода запроса одну строку, внут\-ри которой 


подразумевается связка~<<И>> между терминами. Такой подход приводит 


к~получению значительного <<шума>>, из которого пользователю 


приходится выбирать нужную ему информацию, за\-тра\-чи\-вая значительное 


время, или уходить с~сай\-та биб\-лио\-те\-ки и~больше туда не возвращаться. 


В~качестве иллюстрации такого подхода можно привести каталог 


Российской государственной библиотеки (РГБ), представленный на сайте 


{\sf https://www.rsl.ru}. При попытке узнать, есть ли в~фондах РГБ журнал 


<<Программные продукты и~сис\-те\-мы>> за 2017~г., пользователь, введя 


в~строку запроса <<журнал программные продукты и~сис\-те\-мы 2017~год>> 


получит~12\,830~ссылок на различные издания, в~которых в~разных 


комбинациях встречаются эти слова (рис.~3).


  


\begin{figure} %fig3


\vspace*{1pt}


 \begin{center}


 \mbox{%


 \epsfxsize=128mm 


 \epsfbox{kal-3.eps}


 }


\end{center}


\vspace*{-9pt}


\Caption{Поисковый запрос в~электронном каталоге РГБ}


\end{figure}





  Очевидно, что искать в~этом неупорядоченном массиве нужный журнал ни 


один пользователь не станет. Положение не спасает и~появляющаяся после 


первой выдачи результатов поиска левая часть экрана (см.\ рис.~3), в~которой 


можно указать год издания и~отметить специальность ВАК (Технические 


науки)~--- число выданных документов уменьшается до~611, но в~обозримом 


начале списка нужный журнал отсутствует. В~то же время РГБ должна иметь 


этот журнал, поскольку получает обязательный экземпляр всех изданий, 


выходящих в~Рос\-сии.


{\looseness=1





}





  


  Подобные библиотечные каталоги дискредитируют библиотеки, не 


предоставляя пользователям ничего нового по сравнению с~поисковыми 


машинами Интернета.


  


  Автор считает, что интернет-каталог научной библиотеки должен: 


(а)~в~отличие от общих поисковых машин (Yandex, Google, Rambler и~т.\,п.), 


выдавать ВСЮ ту и~ТОЛЬКО ту информацию, которая соответствует 


запросу, сформулированному пользователем; (б)~обладать свойствами 


развитой информационно-поисковой системы, обеспечивающей обработку 


сложных запросов, включающих операторы булевой логики; 


(в)~предоставлять возможность поиска материалов в~фонде библиотеки как 


по всем элементам расширенного библиографического описания, так и~по 


тематике изданий; (г)~обладать прозрачным пользовательским интерфейсом, 


не требующим знания библиотечной терминологии. 


  


  Именно исходя из этих принципов специалистами БЕН РАН 


разрабатываются интернет-каталоги Библиотеки. Развернутую информацию 


о~поисковых возможностях сводного каталога книг и~продолжающихся 


изданий БЕН РАН предполагается опубликовать в~одном из ближайших 


номеров журнала.


  


{\small\frenchspacing


{%\baselineskip=10.8pt


\addcontentsline{toc}{section}{Литература}


\begin{thebibliography}{9}


\bibitem{1-kal}


ГОСТ 7.1-2003 Библиографическая запись. Библиографическое описание. {\sf 


http:// diss.rsl.ru/datadocs/doc\_291wu.pdf.}


\bibitem{2-kal}


\Au{Каленов Н.\,Е.} Электронные каталоги академических библиотек: какими им 


быть?~/\!\!/ Теория и~практика об\-ще\-ст\-вен\-но-на\-уч\-ной информации, 2014. Вып.~22. 


С.~54--63.


\bibitem{3-kal}


\Au{Власова С.\,А., Каленов Н.\,Е.} Роль каталогов научных библиотек в~задачах 


информационного сопровождения научных исследований~/\!\!/ Информационные 


процессы, 2014. Т.~14. №\,3. С.~232--241. {\sf http://www.jip.ru/2014/232-241-2014.htm}.


\bibitem{4-kal}


\Au{Каленов Н.\,Е.} Библиотечные интернет-каталоги и~пользователи~/\!\!/ Научная 


периодика: проблемы и~решения, 2015. Т.~5. №\,6. С.~265--272.


        \end{thebibliography}


} }








\vspace*{-3pt}





\hfill{\small\textit{Поступила в~редакцию 24.04.18}}














\vspace*{8pt}





\hrule





\vspace*{2pt}





%\newpage





%\vspace*{-28pt}





\hrule





%\vspace*{9pt}











\def\tit{SOME EVALUATION OF~THE~RUSSIAN LIBRARIES ONLINE 


CATALOGS SEARCH CAPABILITIES}





\def\titkol{Some evaluation of~the~Russian libraries online 


catalogs search capabilities}








\def\aut{N.\,E.~Kalenov}


\def\autkol{N.\,E.~Kalenov}











\titel{\tit}{\aut}{\autkol}{\titkol}





\vspace*{-9pt}





\noindent


  The Library for Natural Science of the Russian Academy of Sciences, 


11\!/\!11~Znamenka Str., Moscow 119991, Russian Federation








 \noindent





\def\leftfootline{\small{\textbf{\thepage}


\hfill Sistemy i Sredstva Informatiki~---


Systems and Means of Informatics 2018 vol~28 no\,3}


}%


 \def\rightfootline{\small{Sistemy i Sredstva Informatiki~---


 Systems and Means of Informatics   2018   vol~28   no\,3


\hfill \textbf{\thepage}}}





  


  


  


  \Abste{The problems of organizing the information retrieval in library catalogs 


both in the ``historical'' aspect (with reference to traditional card catalogs and 


printed catalogs) and in the modern aspect, conditioned by the development of\linebreak\vspace*{-12pt}}





\Abstend{information technologies, Internet, and universal search engines, are under


consideration. The analysis is made and the shortcomings of the existing online 


catalogs of various scientific libraries of the country are revealed. A number of 


examples of incorrect query processing by catalogs of leading libraries of Russia 


are considered. The conclusion is made about the necessity of a new approach to 


the design of scientific libraries online catalogs. These catalogs should give the 


user all the information he\!/\!she needs without information noise. When creating 


software shells of modern online catalogs, on the one hand, it is necessary to move 


away from the traditional library ideas about catalogs, significantly limiting their 


search capabilities, and on the other hand~--- from the logic of the universal Internet 


search engines, giving the user thousands of percent of information noise.}


  


  \KWE{electronic catalogs; Internet; bibliographic descriptions; subject 


headings; network technologies; search requests; information retrieval}








\DOI{10.14357\!/\!08696527180319} 





%\vspace*{-24pt}





%\Ack


%\noindent





 








\renewcommand{\bibname}{\protect\rmfamily References}


%\renewcommand{\bibname}{\large\protect\rm References}





{\small\frenchspacing


{%\baselineskip=10.8pt


\addcontentsline{toc}{section}{References}


\begin{thebibliography}{9}


\bibitem{1-kal-1}


GOST 7.1-2003 Bibliograficheskaya zapis'. 


Bibliograficheskoe opisanie [Bibliographic record. 


Bibliographic description]. Available at: {\sf 


http://diss.rsl.ru/datadocs/ doc\_291wu.pdf} (accessed August~28, 2018).


\bibitem{2-kal-1}


\Aue{Kalenov, N.\,E.} 2014. Elektronnye katalogi akademicheskikh bibliotek: 


kakimi im byt'? [Electronic catalogs of academic libraries: What should they 


be?]. \textit{Teoriya i~praktika obshchestvenno-nauchnoy informatsii}


[Theory and Practice of Social and Scientific Information] 22:54--63.


\bibitem{3-kal-1}


\Aue{Vlasova, S.\,A., and N.\,E.~Kalenov.} 2014. Rol' katalogov


nauchnykh bibliotek v~za\-da\-chakh informatsionnogo soprovozhdeniya


nauchnykh issledovaniy


[The role of academic 


libraries catalogues for the scientific research information support]. 


\textit{Information Processes} 14(3):232--241. Available 


at: {\sf http://www.jip.ru/2014/232-241-2014.htm}  (accessed August~28, 2018).


\bibitem{4-kal-1}


\Aue{Kalenov, N.\,E.} 2015. Bibliotechnye internet-katalogi i~pol'zovateli


[Online library catalogs and users]. 


\textit{Scientific Periodicals: Problems and Solutions} 5(6):265--272.


\end{thebibliography}


} }








\vspace*{-3pt}





\hfill{\small\textit{Received April 24, 2018}}





\Contrl





\noindent


\textbf{Kalenov Nikolay E.} (b.\ 1945)~--- Doctor of Science in 


technology, professor, principal scientist, Library for Natural 


Sciences of the Russian Academy of Sciences, 11\!/\!11~Znamenka Str., 


Moscow 119991, Russian Federation; \mbox{nek@benran.ru}











\label{end\stat}








\renewcommand{\bibname}{\protect\rm Литература}


